% !TEX root = ./slide_main.tex
% 複数フォルダに分かれている場合、コンパイルするディレクトリがルートになる場合があるため必要
%==============================================================================
% 卒業論文・修士論文 発表用 Beamer スライドテンプレート
%
% コンパイル方法（このフォルダで実行）:
%   latexmk slide_main.tex     →   slide_main.pdf が生成される
%
% 編集手順:
%   1. 下の「タイトル情報」を自分のものに書き換える
%   2. 各スライドのサンプルを参考に、自分の研究内容に置き換える
%   3. 不要なスライドは削除、足りないスライドはコピーして増やす
%==============================================================================
\documentclass[uplatex,dvipdfmx,aspectratio=169]{beamer}
% aspectratio=169 : 16:9 ワイド画面（4:3 にしたい場合は aspectratio=43 に変更）
%
% 配布用PDF（handout出力）にする場合:
%   \pause を無視して 1フレーム=1ページ にまとめたPDFを生成する。
%   → 下の行に差し替えてコンパイルする（README参照）
%   \documentclass[uplatex,dvipdfmx,aspectratio=169,handout]{beamer}
% 配布用・pause無視（1フレーム=1ページ）:
% \documentclass[uplatex,dvipdfmx,aspectratio=169,handout]{beamer}

%==============================================================================
% 卒論・修論発表用 Beamer スライドテンプレート - ヘッダーファイル
%
% このファイルにはスライドで使用する各種設定が含まれています。
% slide_main.tex から %==============================================================================
% 卒論・修論発表用 Beamer スライドテンプレート - ヘッダーファイル
%
% このファイルにはスライドで使用する各種設定が含まれています。
% slide_main.tex から %==============================================================================
% 卒論・修論発表用 Beamer スライドテンプレート - ヘッダーファイル
%
% このファイルにはスライドで使用する各種設定が含まれています。
% slide_main.tex から \input{slide_preamble} で読み込まれます。
% テーマはデフォルト（\usetheme の指定なし）です。
% テーマを変えたい場合は slide_main.tex で \usetheme{...} を指定してください。
%==============================================================================

%------------------ 日本語フォント設定 ----------------------------------------
% 和文をゴシック体（サンセリフ）に統一する。
% Beamer は欧文がサンセリフ体のため、和文もゴシックに合わせると見た目が揃う。
\renewcommand{\kanjifamilydefault}{\gtdefault}

%------------------ ハイパーリンク ---------------------------------------------
% PDF のしおり（ブックマーク）を日本語対応にする（uplatex + dvipdfmx 用）。
% beamer は内部で hyperref を読み込むため、\documentclass の直後に読み込む。
\usepackage{pxjahyper}

%------------------ 数式 --------------------------------------------------------
% beamer は amsmath, amssymb を自動で読み込むため、追加は不要。
% （必要なら以下をコメントアウト解除して追加する）
% \usepackage{mathtools}   % \coloneqq など高度な数式マクロ

%------------------ 画像・表 ----------------------------------------------------
\usepackage{graphicx}    % 画像挿入（\includegraphics）
\usepackage{booktabs}    % 表の罫線（\toprule / \midrule / \bottomrule）
\usepackage{subcaption}  % 複数図を (a), (b) 形式でまとめる（fig:xx-a 等）

%------------------ 動画 --------------------------------------------------------
% 動画を PDF に埋め込む機能は「使用不可」。
% 理由: media9 は Flash Player（終了済み）ベースのため実質再生できないうえ、
%       PDF が約13MBに肥大化する。
% \usepackage{media9} はコメントアウト済み（下の行）。
% どうしても使う場合は slide_main.tex の動画スライド（コメントアウト中）も有効化する。
% \usepackage{media9} % Flash Player死亡のため使用不可
% 使用例:
%   \includemedia[width=0.55\linewidth,height=0.31\linewidth,activate=onclick,
%                 flashvars={autoplay=true}]{ポスター画像}{movie/動画ファイル.mp4}
% 注意: multimedia パッケージの \movie は dvipdfmx と非互換のため使用しないこと

%------------------ 定理環境 ----------------------------------------------------
% beamer には block / alertblock / exampleblock と、
% theorem, lemma, corollary, definition, example, proof が標準で用意されている。
% （スライドでは \begin{block}{タイトル} を主に使うとよい。
%   タイトルは日本語で指定できるので \newtheorem は定義していない。）

%------------------ よく使う数式コマンド ---------------------------------------
\newcommand{\bm}{\boldsymbol}       % 太字ベクトル: \bm{x}
\newcommand{\del}{\partial}         % 偏微分: \del
\newcommand{\diag}{\mathrm{diag}}   % 対角行列: \diag

%------------------ フッター（ページ番号） --------------------------------------
% 右下のナビゲーション記号（ゴミのようなアイコン）を非表示
\setbeamertemplate{navigation symbols}{}
% フッター: 右下に「ページ番号 / 総ページ数」を表示
% （タイトルページや [plain] 指定のフレームには表示されない）
\setbeamertemplate{footline}{%
  \hbox{%
    \begin{beamercolorbox}[wd=\paperwidth,ht=2.5ex,dp=1.125ex]{}%
      \hfill\insertframenumber\,/\,\inserttotalframenumber\hspace{2mm}%
    \end{beamercolorbox}%
  }%
}

%------------------ セクションページ -------------------------------------------
% 各セクションの開始時に、セクションタイトルのスライドを自動挿入する
% 注意: \section*{...} でも発火する（目次には載らないが区切りページは挿入される）
\AtBeginSection{%
  \begin{frame}[plain]
    \sectionpage
  \end{frame}%
}
 で読み込まれます。
% テーマはデフォルト（\usetheme の指定なし）です。
% テーマを変えたい場合は slide_main.tex で \usetheme{...} を指定してください。
%==============================================================================

%------------------ 日本語フォント設定 ----------------------------------------
% 和文をゴシック体（サンセリフ）に統一する。
% Beamer は欧文がサンセリフ体のため、和文もゴシックに合わせると見た目が揃う。
\renewcommand{\kanjifamilydefault}{\gtdefault}

%------------------ ハイパーリンク ---------------------------------------------
% PDF のしおり（ブックマーク）を日本語対応にする（uplatex + dvipdfmx 用）。
% beamer は内部で hyperref を読み込むため、\documentclass の直後に読み込む。
\usepackage{pxjahyper}

%------------------ 数式 --------------------------------------------------------
% beamer は amsmath, amssymb を自動で読み込むため、追加は不要。
% （必要なら以下をコメントアウト解除して追加する）
% \usepackage{mathtools}   % \coloneqq など高度な数式マクロ

%------------------ 画像・表 ----------------------------------------------------
\usepackage{graphicx}    % 画像挿入（\includegraphics）
\usepackage{booktabs}    % 表の罫線（\toprule / \midrule / \bottomrule）
\usepackage{subcaption}  % 複数図を (a), (b) 形式でまとめる（fig:xx-a 等）

%------------------ 動画 --------------------------------------------------------
% 動画を PDF に埋め込む機能は「使用不可」。
% 理由: media9 は Flash Player（終了済み）ベースのため実質再生できないうえ、
%       PDF が約13MBに肥大化する。
% \usepackage{media9} はコメントアウト済み（下の行）。
% どうしても使う場合は slide_main.tex の動画スライド（コメントアウト中）も有効化する。
% \usepackage{media9} % Flash Player死亡のため使用不可
% 使用例:
%   \includemedia[width=0.55\linewidth,height=0.31\linewidth,activate=onclick,
%                 flashvars={autoplay=true}]{ポスター画像}{movie/動画ファイル.mp4}
% 注意: multimedia パッケージの \movie は dvipdfmx と非互換のため使用しないこと

%------------------ 定理環境 ----------------------------------------------------
% beamer には block / alertblock / exampleblock と、
% theorem, lemma, corollary, definition, example, proof が標準で用意されている。
% （スライドでは \begin{block}{タイトル} を主に使うとよい。
%   タイトルは日本語で指定できるので \newtheorem は定義していない。）

%------------------ よく使う数式コマンド ---------------------------------------
\newcommand{\bm}{\boldsymbol}       % 太字ベクトル: \bm{x}
\newcommand{\del}{\partial}         % 偏微分: \del
\newcommand{\diag}{\mathrm{diag}}   % 対角行列: \diag

%------------------ フッター（ページ番号） --------------------------------------
% 右下のナビゲーション記号（ゴミのようなアイコン）を非表示
\setbeamertemplate{navigation symbols}{}
% フッター: 右下に「ページ番号 / 総ページ数」を表示
% （タイトルページや [plain] 指定のフレームには表示されない）
\setbeamertemplate{footline}{%
  \hbox{%
    \begin{beamercolorbox}[wd=\paperwidth,ht=2.5ex,dp=1.125ex]{}%
      \hfill\insertframenumber\,/\,\inserttotalframenumber\hspace{2mm}%
    \end{beamercolorbox}%
  }%
}

%------------------ セクションページ -------------------------------------------
% 各セクションの開始時に、セクションタイトルのスライドを自動挿入する
% 注意: \section*{...} でも発火する（目次には載らないが区切りページは挿入される）
\AtBeginSection{%
  \begin{frame}[plain]
    \sectionpage
  \end{frame}%
}
 で読み込まれます。
% テーマはデフォルト（\usetheme の指定なし）です。
% テーマを変えたい場合は slide_main.tex で \usetheme{...} を指定してください。
%==============================================================================

%------------------ 日本語フォント設定 ----------------------------------------
% 和文をゴシック体（サンセリフ）に統一する。
% Beamer は欧文がサンセリフ体のため、和文もゴシックに合わせると見た目が揃う。
\renewcommand{\kanjifamilydefault}{\gtdefault}

%------------------ ハイパーリンク ---------------------------------------------
% PDF のしおり（ブックマーク）を日本語対応にする（uplatex + dvipdfmx 用）。
% beamer は内部で hyperref を読み込むため、\documentclass の直後に読み込む。
\usepackage{pxjahyper}

%------------------ 数式 --------------------------------------------------------
% beamer は amsmath, amssymb を自動で読み込むため、追加は不要。
% （必要なら以下をコメントアウト解除して追加する）
% \usepackage{mathtools}   % \coloneqq など高度な数式マクロ

%------------------ 画像・表 ----------------------------------------------------
\usepackage{graphicx}    % 画像挿入（\includegraphics）
\usepackage{booktabs}    % 表の罫線（\toprule / \midrule / \bottomrule）
\usepackage{subcaption}  % 複数図を (a), (b) 形式でまとめる（fig:xx-a 等）

%------------------ 動画 --------------------------------------------------------
% 動画を PDF に埋め込む機能は「使用不可」。
% 理由: media9 は Flash Player（終了済み）ベースのため実質再生できないうえ、
%       PDF が約13MBに肥大化する。
% \usepackage{media9} はコメントアウト済み（下の行）。
% どうしても使う場合は slide_main.tex の動画スライド（コメントアウト中）も有効化する。
% \usepackage{media9} % Flash Player死亡のため使用不可
% 使用例:
%   \includemedia[width=0.55\linewidth,height=0.31\linewidth,activate=onclick,
%                 flashvars={autoplay=true}]{ポスター画像}{movie/動画ファイル.mp4}
% 注意: multimedia パッケージの \movie は dvipdfmx と非互換のため使用しないこと

%------------------ 定理環境 ----------------------------------------------------
% beamer には block / alertblock / exampleblock と、
% theorem, lemma, corollary, definition, example, proof が標準で用意されている。
% （スライドでは \begin{block}{タイトル} を主に使うとよい。
%   タイトルは日本語で指定できるので \newtheorem は定義していない。）

%------------------ よく使う数式コマンド ---------------------------------------
\newcommand{\bm}{\boldsymbol}       % 太字ベクトル: \bm{x}
\newcommand{\del}{\partial}         % 偏微分: \del
\newcommand{\diag}{\mathrm{diag}}   % 対角行列: \diag

%------------------ フッター（ページ番号） --------------------------------------
% 右下のナビゲーション記号（ゴミのようなアイコン）を非表示
\setbeamertemplate{navigation symbols}{}
% フッター: 右下に「ページ番号 / 総ページ数」を表示
% （タイトルページや [plain] 指定のフレームには表示されない）
\setbeamertemplate{footline}{%
  \hbox{%
    \begin{beamercolorbox}[wd=\paperwidth,ht=2.5ex,dp=1.125ex]{}%
      \hfill\insertframenumber\,/\,\inserttotalframenumber\hspace{2mm}%
    \end{beamercolorbox}%
  }%
}

%------------------ セクションページ -------------------------------------------
% 各セクションの開始時に、セクションタイトルのスライドを自動挿入する
% 注意: \section*{...} でも発火する（目次には載らないが区切りページは挿入される）
\AtBeginSection{%
  \begin{frame}[plain]
    \sectionpage
  \end{frame}%
}


% テーマを変えたい場合はここで指定する。
% 例: \usetheme{Madrid}, \usetheme{CambridgeUS} など
% 注意: metropolis など XeLaTeX/LuaLaTeX 必須のテーマは uplatex では使えない
\usetheme{Madrid}

%==============================================================================
% タイトル情報（ここだけ各自で編集）
%==============================================================================
\title[卒論・修論発表]{卒業論文・修士論文のタイトル}
\subtitle{サブタイトル（任意・省略可）}
\author{氏名を記入}
\institute{福井大学 工学部 機械・システム工学科 機械工学コース}
\date{\number\year 年 \number\month 月 \number\day 日}  % 例: 2026年9月1日
% タイトルページ下部に表示する情報（指導教員など）
\titlegraphic{%
  {\small 指導教員：〇〇 教授}%
}

\begin{document}

%==============================================================================
% タイトルページ
%==============================================================================
\begin{frame}[plain]
  \titlepage
\end{frame}

%==============================================================================
% 目次
%==============================================================================
\begin{frame}{目次}
  \tableofcontents
  % 目次に載せたくないセクションは \section*{...} を使う（例: 参考文献・謝辞）
\end{frame}

%==============================================================================
% はじめに（背景・目的）
%==============================================================================
\section{はじめに}

\subsection{研究背景}
\begin{frame}{研究背景}
  \begin{itemize}
    \item 研究の動機や社会的・学術的な背景を簡潔に書く
    \item 図や表があると伝わりやすい
    \item 参考文献を引用する場合は論文と同じように \cite{ref1} を使う
  \end{itemize}
  \begin{block}{背景のポイント}
    \begin{itemize}
      \item 「なぜこの研究をするのか」を聴衆に伝える
      \item 1スライドに収めるのが目安
    \end{itemize}
  \end{block}
\end{frame}

\subsection{研究目的}
\begin{frame}{研究目的}
  \begin{enumerate}
    \item \textbf{目的1}：解決したい課題を具体的に書く
    \item \textbf{目的2}：達成したいことを定量的に書く
    \item \textbf{目的3}：対象・範囲・条件を明示する
  \end{enumerate}
  \pause
  \begin{exampleblock}{本発表の流れ}
    背景 → 関連研究 → 提案手法 → 実験・結果 → 考察 → まとめ
  \end{exampleblock}
\end{frame}

%==============================================================================
% 関連研究
%==============================================================================
\section{関連研究}

\begin{frame}{既存手法と課題}
  \begin{columns}[t]
    \begin{column}{0.48\linewidth}
      \textbf{既存手法}
      \begin{itemize}
        \item 手法Aの長所・短所
        \item 手法Bの長所・短所
      \end{itemize}
    \end{column}
    \begin{column}{0.48\linewidth}
      \textbf{課題}
      \begin{itemize}
        \item 未解決の問題1
        \item 未解決の問題2
      \end{itemize}
    \end{column}
  \end{columns}

  \vspace{3mm}
  \begin{table}
    \centering
    \caption{既存手法との比較（記入例）}
    \label{tab:comparison}
    \begin{tabular}{lccc}
      \toprule
      & 手法A & 手法B & 提案手法 \\
      \midrule
      精度 & 高い & 低い & \textbf{高い} \\
      計算コスト & 大 & 小 & \textbf{小} \\
      実装の容易さ & 難 & 易 & \textbf{易} \\
      \bottomrule
    \end{tabular}
  \end{table}
\end{frame}
%==============================================================================
% 提案手法
%==============================================================================
\section{提案手法}

\begin{frame}{提案手法の概要}
  \begin{block}{提案手法のアイデア}
    手法の概要を 1〜2 文で簡潔に。
  \end{block}
  \begin{enumerate}
    \item ステップ1：前処理（データの説明）
    \item ステップ2：特徴抽出（数式は次スライド参照）
    \item ステップ3：学習・判定（アルゴリズムの説明）
  \end{enumerate}
  \pause
  \begin{alertblock}{注意点}
    既存手法との違い・新規性をここで明確に示す。
  \end{alertblock}
\end{frame}

\begin{frame}{定式化}
  提案手法の数式の例:
  \begin{equation}
    \label{eq:objective}
    \min_{\bm{w}} \; \frac{1}{2}\|\bm{w}\|^2 + C\sum_{i=1}^{n}\xi_i
  \end{equation}
  \begin{itemize}
    \item $\bm{w}$：重みベクトル
    \item $\xi_i$：スラック変数（$i=1,\dots,n$）
    \item $C$：正則化パラメータ
  \end{itemize}
  \pause
  複数行の式は align 環境を使う:
  \begin{align}
    y &= f(\bm{x}) = \bm{w}^{\top}\bm{x} + b \\
    \hat{y} &= \mathrm{sign}(y)
  \end{align}
\end{frame}

%==============================================================================
% 実験・結果
%==============================================================================
\section{実験・結果}

\begin{frame}{実験条件}
  \begin{table}
    \centering
    \caption{実験条件（記入例）}
    \label{tab:condition}
    \begin{tabular}{ll}
      \toprule
      項目 & 内容 \\
      \midrule
      データセット & 実験データ名（サンプル数） \\
      比較手法 & 既存手法A、既存手法B \\
      評価指標 & 精度（Accuracy）、F値 など \\
      環境 & CPU / GPU、フレームワーク名、言語 \\
      \bottomrule
    \end{tabular}
  \end{table}
  \pause
  実験条件は再現可能なように具体的に書く。
\end{frame}

\begin{frame}{実験結果（図）}
  \begin{figure}
    \centering
    \includegraphics[width=0.55\linewidth]{fig/fig_1.png}
    \caption{実験結果の例（fig/fig\_1.png）}
    % 注意: 本文中のアンダースコアは \_ と書くこと（fig_1.png と書くとエラー）
    \label{fig:result1}
  \end{figure}
\end{frame}

\begin{frame}{実験結果（比較）}
  \begin{columns}[c]
    \begin{column}{0.48\linewidth}
      \begin{figure}
        \centering
        \includegraphics[width=\linewidth]{fig/fig_1.png}
        \caption{既存手法}
        \label{fig:result2-a}
      \end{figure}
    \end{column}
    \begin{column}{0.48\linewidth}
      \begin{figure}
        \centering
        \includegraphics[width=\linewidth]{fig/fig_2.png}
        \caption{提案手法}
        \label{fig:result2-b}
      \end{figure}
    \end{column}
  \end{columns}
\end{frame}

% subcaption パッケージを使った複数図の例（使う場合はコメントを外す）:
% \begin{frame}{実験結果（subcaption の例）}
%   \begin{figure}
%     \centering
%     \begin{subfigure}{0.48\linewidth}
%       \centering\includegraphics[width=\linewidth]{fig/fig_1.png}
%       \subcaption{既存手法}\label{fig:result3-a}
%     \end{subfigure}
%     \hfill
%     \begin{subfigure}{0.48\linewidth}
%       \centering\includegraphics[width=\linewidth]{fig/fig_2.png}
%       \subcaption{提案手法}\label{fig:result3-b}
%     \end{subfigure}
%     \caption{結果の比較}\label{fig:result3}
%   \end{figure}
% \end{frame}

%==============================================================================
% 実験・結果（動画）
%==============================================================================
% 動画はあきらめたほうがよい
% \begin{frame}{実験結果（動画）}
%   \begin{center}
%     \includemedia[
%       width=0.55\linewidth,
%       height=0.31\linewidth,
%       activate=onclick,
%       flashvars={autoplay=true}
%     ]{\fcolorbox{black}{black}{\color{white}\Large\bfseries ▶ クリックして動画を再生}}%
%     {movie/soft_material_recognition.mp4}
%   \end{center}
%   \begin{center}
%     {\small 実験の様子（クリックで再生。Adobe Acrobat / Acrobat Reader が必要）}
%   \end{center}
% \end{frame}

%==============================================================================
% 考察・まとめ
%==============================================================================
\section{考察とまとめ}

\begin{frame}{考察}
  \begin{itemize}
    \item 結果から言えることを述べる
    \item 既存手法と比較してどう優れているか / 劣っているか
    \item 結果が期待通りでなかった場合はその理由を考察する
  \end{itemize}
  \pause
  \begin{alertblock}{注意}
    事実（結果）と意見（考察）を明確に区別して話すこと。
  \end{alertblock}
\end{frame}

\begin{frame}{まとめと今後の課題}
  \begin{columns}[t]
    \begin{column}{0.48\linewidth}
      \textbf{まとめ}
      \begin{itemize}
        \item 目的に対して何が達成できたか
        \item 成果の要点を3つ程度に絞る
      \end{itemize}
    \end{column}
    \begin{column}{0.48\linewidth}
      \textbf{今後の課題}
      \begin{itemize}
        \item 残っている課題1
        \item 課題2（改善の見通し）
      \end{itemize}
    \end{column}
  \end{columns}
\end{frame}

%==============================================================================
% 参考文献（目次に載せない場合: \section* を使う）
% 注意: \section* でもセクション区切りページは挿入される。
% 目次に載せて区切りページも見せたい場合は \section{参考文献} にする。
%==============================================================================
\section*{参考文献}
\begin{frame}{参考文献}
  \begin{thebibliography}{99}
    \bibitem{ref1}
      著者名, ``論文タイトル,'' 学会名・ジャーナル名, 巻, 号, pp.xx-xx, 年.
    \bibitem{ref2}
      著者名, ``論文タイトル,'' 学会名・ジャーナル名, 巻, 号, pp.xx-xx, 年.
  \end{thebibliography}
  % bibtex を使う場合は、上記の代わりに
  % \bibliographystyle{ssice}
  % \bibliography{bibfile}
  % と書いて bibfile.bib を用意する（他テンプレートと同様）。
\end{frame}

%==============================================================================
% 謝辞・おわり
%==============================================================================
\begin{frame}[plain]
  \centering
  \vspace{2cm}
  {\Huge\bfseries ご清聴ありがとうございました}\\[1.5em]
  {\Large 質疑応答をお願いします}\\[3cm]
  {\small 連絡先や謝辞が必要ならここに書く}
\end{frame}

\end{document}

